\chapter{État de l'art}
\label{chap:etat_art}

\section*{Introduction}
La compréhension automatique des émotions humaines dans le texte constitue l'un des défis centraux du Traitement Automatique du Langage Naturel (TALN). Longtemps abordée comme un problème de classification mono-étiquette visant à identifier l'émotion dominante d'un énoncé, la détection des émotions connaît aujourd'hui un tournant paradigmatique : la réalité du langage humain révèle que plusieurs émotions peuvent coexister simultanément dans un même texte, s'interpénétrer, se moduler et interagir à différentes granularités linguistiques.

\noindent Le présent travail s'inscrit dans ce contexte en proposant l'architecture \textbf{HAKE-MER} (\textit{Hierarchical Attention with Knowledge Enhancement for Multi-Emotion Recognition}), une approche qui ambitionne de dépasser les limitations des méthodes existantes en combinant : \\
(i) une modélisation hiérarchique multi-niveaux basée sur les transformers \cite{vaswani2017,devlin2019,yang2016}, \\
(ii) un mécanisme d'attention croisée pour capturer les interactions inter-émotions,\\
(iii) une intégration de connaissances sémantiques externes issues de bases telles que \textit{ConceptNet} \cite{speer2017} et \textit{EmoNet}.

\noindent Le présent état de l'art explore les cinq axes thématiques suivantes : les méthodes de détection d'émotions multiples, les architectures d'attention hiérarchique et les transformers, l'intégration de connaissances externes, les techniques de raisonnement par bon sens, et les approches d'apprentissage multi-niveaux et multi-granularités.

\section{Méthodes de Détection d'Émotions Multiples (Multi-label Emotion Classification)}

\subsection{Fondements et cadrage du problème}
\noindent La détection des émotions dans le texte s'appuie historiquement sur deux théories fondatrices : la théorie des émotions de base d'Ekman (1992), qui postule l'existence de six émotions universelles (joie, tristesse, peur, colère, dégoût, surprise), et la théorie dimensionnelle de Russell (1980), fondée sur un espace bidimensionnel valence-arousal (V-A). Ces deux approches ont structuré les premiers systèmes de classification, mais leurs limitations sont manifestes dès lors que l'on confronte la richesse émotionnelle du langage naturel à ces catégories rigides.

\noindent La classification multi-étiquettes (\textit{multi-label classification}, MLC) constitue le cadre formel adapté à la détection simultanée de plusieurs émotions. Formellement, étant donné un énoncé $x$ et un ensemble d'étiquettes émotionnelles $L = \{l_{1}, l_{2}, \dots, l_{n}\}$, la tâche consiste à apprendre une fonction
\[
f : x \rightarrow \{0,1\}^{n}
\]
assignant à chaque étiquette une valeur binaire ou un score de probabilité. Cette formulation diffère fondamentalement de la classification multi-classes, où une seule étiquette est attribuée.

\subsection{Approches basées sur les features classiques}
\noindent Les premières approches de classification émotionnelle multi-étiquettes reposaient sur l'extraction manuelle de traits linguistiques : n-grammes, présence de mots affectifs issus de lexiques (\textit{LIWC}, \textit{SentiWordNet}, \textit{NRC Emotion Lexicon}), structures syntaxiques et traits prosodiques. Ces représentations étaient ensuite classifiées par des algorithmes tels que les \textit{Support Vector Machines} (SVM) en configuration multi-label, les forêts aléatoires, ou les chaînes de classificateurs (\textit{Classifier Chains}) de Read et al. (2009), qui modélisent explicitement les dépendances entre étiquettes.

\noindent La limite majeure de ces méthodes réside dans leur incapacité à capturer la sémantique profonde et les dépendances contextuelles à longue portée, ainsi que leur fragilité face aux expressions figuratives, à l'ironie et aux émotions implicites.

\subsection{Approches par apprentissage profond}
\noindent L'avènement du deep learning a profondément reconfiguré le paysage de la détection émotionnelle multi-étiquettes. Les réseaux convolutionnels (CNN) de Kim (2014) ont démontré l'efficacité des filtres locaux pour la classification de textes courts. Les réseaux récurrents bidirectionnels (BiLSTM) ont ensuite permis de capturer les dépendances séquentielles dans les deux sens de lecture, surpassant les CNN sur des tâches nécessitant une compréhension contextuelle étendue.

\noindent Un apport méthodologique majeur dans le contexte multi-label est le mécanisme d'attention de Bahdanau et al. (2015) \cite{bahdanau2015}, qui permet au modèle de pondérer différemment les tokens selon leur pertinence pour chaque étiquette. Des extensions telles que l'attention spécifique à une étiquette (\textit{label-specific attention}) ont été proposées pour modéliser individuellement la contribution de chaque région textuelle à chaque émotion.

\subsection{Défis spécifiques : labellisation incomplète et subjectivité}
\noindent La construction de datasets pour la détection multi-émotions est confrontée à deux problèmes structurels majeurs. D'une part, la labellisation incomplète (\textit{partial label learning}) : les annotateurs tendent à ne marquer que les émotions saillantes, omettant les émotions de faible intensité ou secondaires. D'autre part, la subjectivité inhérente à la perception émotionnelle fait que différents annotateurs peuvent attribuer des étiquettes contradictoires au même énoncé.

\noindent Des approches de \textit{learning with noisy labels}, d'agrégation bayésienne des annotations (Dawid \& Skene, 1979 ; Rodrigues \& Pereira, 2018) et de \textit{curriculum learning} ont été proposées pour atténuer ces biais. La modélisation probabiliste des distributions d'étiquettes plutôt que des affectations binaires constitue une direction prometteuse explorée notamment par Demszky et al. (2020) dans la conception du dataset \textit{GoEmotions} \cite{demszky2020}.

\section{Architectures d'Attention Hiérarchique et Transformers}

\subsection{Le mécanisme d'attention : fondements théoriques}
\noindent Le mécanisme d'attention, introduit formellement par Bahdanau et al. (2015) \cite{bahdanau2015} dans le cadre des modèles séquence-à-séquence, puis généralisé sous la forme d'une attention multi-têtes (\textit{multi-head attention}) par Vaswani et al. (2017) \cite{vaswani2017} dans l'architecture \textit{Transformer}, constitue le socle théorique des approches modernes de TALN. L'attention permet à un modèle de calculer, pour chaque position d'une séquence, une pondération sur l'ensemble des autres positions, capturant ainsi des dépendances arbitrairement longues sans les limitations des architectures récurrentes.

Formellement, l'attention est calculée comme :
\[
\text{Attention}(Q,K,V)=\text{softmax}\left(\frac{QK^{T}}{\sqrt{d_k}}\right)V
\]
où $Q$, $K$, $V$ désignent respectivement les matrices \textit{query}, \textit{key} et \textit{value}, et $d_k$ est la dimension des clés. Cette formulation est au cœur de l'architecture \textit{Transformer} et de tous ses descendants \cite{vaswani2017}.

\subsection{L'architecture Transformer et ses variants pré-entraînés}
\noindent L'architecture \textit{Transformer} \cite{vaswani2017} a engendré une famille de modèles pré-entraînés sur de larges corpus textuels, qui ont révolutionné pratiquement toutes les tâches de TALN. BERT \cite{devlin2019} introduit un pré-entraînement bidirectionnel par masquage de tokens (\textit{Masked Language Modeling}) et prédiction de phrase suivante, produisant des représentations contextuelles riches. RoBERTa (Liu et al., 2019) améliore BERT par un entraînement plus long sur un corpus plus large et la suppression de la tâche NSP. XLNet (Yang et al., 2019) adopte un schéma autorégressif permutationnel contournant les limitations du masquage.

\noindent Pour le domaine affectif, des variantes spécialisées ont été développées : RoBERTa-base fine-tuné sur des datasets émotionnels \cite{demszky2020}, SentiBERT (Yin et al., 2020) enrichi de structures syntaxiques, et EmotionBERT. Ces modèles constituent les encodeurs de base des approches \textit{state-of-the-art} pour la classification émotionnelle multi-étiquettes.

\subsection{Architectures d'attention hiérarchique}
\noindent La modélisation hiérarchique du texte vise à capturer les structures à différentes granularités : mot, phrase (ou clause), et document. Le \textit{Hierarchical Attention Network} (HAN) proposé par Yang et al. (2016) \cite{yang2016} pour la classification de documents constitue le modèle fondateur de cette famille d'approches. HAN empile deux niveaux d'encodeurs BiGRU avec attention : un au niveau des mots au sein de chaque phrase, et un au niveau des phrases au sein du document, permettant de capturer la structure rhétorique du texte.

\noindent Des extensions hiérarchiques des transformers ont ensuite été proposées pour traiter des textes longs. Le modèle HIBERT (Zhang et al., 2019) applique des encodeurs transformer à deux niveaux. Des architectures récentes comme HAT (\textit{Hierarchical Attention Transformer}) utilisent des self-attentions locales au niveau des segments et une attention globale au niveau du document, réduisant la complexité computationnelle tout en maintenant une modélisation multi-échelle.

\noindent Dans le contexte émotionnel spécifique, Luo et al. (2025) proposent un mécanisme d'attention à trois dimensions pour la classification multi-étiquettes des émotions : une couche d'attention distincte par émotion, chacune traitée comme un plan bidimensionnel, reliées entre elles par une attention multi-têtes, avec un paramètre de polarité/intensité (\textit{Commander}) modulant les poids d'attention. Évaluée sur \textit{GoEmotions} et \textit{SemEval-2018}, cette approche porte le F1-macro de 0,46 à 0,56 sur GoEmotions \cite{luo2025attention}. Dans un cadre plus général de classification multi-étiquettes, Li et al. (2022) proposent une architecture Transformer-CNN hiérarchique combinant une attention multi-têtes au niveau des mots et un CNN au niveau des phrases pour capturer l'information sémantique à plusieurs échelles \cite{li2022hierarchical} ; bien que non spécifique aux émotions, ce principe de traitement hiérarchique multi-échelle rejoint directement la motivation du module d'encodage hiérarchique de HAKE-MER.

\subsection{Attention croisée pour les interactions entre émotions}
\noindent La modélisation des interactions entre émotions constitue un verrou scientifique central. Les émotions ne sont pas indépendantes : la joie et la surprise sont fréquemment co-occurrentes, tandis que la joie et la tristesse s'excluent plus souvent. Des mécanismes d'attention croisée (\textit{cross-attention}) ou d'attention inter-étiquettes permettent de capturer ces corrélations. Des approches comme le \textit{Label-Correlation Attention} de Chen et al. (2023) construisent un graphe d'émotions et appliquent des réseaux d'attention sur ce graphe pour propager les signaux entre étiquettes corrélées.

\section{Intégration de Connaissances Externes}

\subsection{Motivation et cadre général}
\noindent L'intégration de connaissances externes dans les modèles de TALN répond à une limitation fondamentale des approches purement statistiques : la difficulté à inférer les émotions implicites, à comprendre les références culturelles ou à désambiguïser les expressions indirectes. Les connaissances structurées permettent d'ancrer les représentations dans une sémantique explicite et d'effectuer des inférences au-delà de la co-occurrence statistique.

\subsection{Bases de connaissances affectives}
\noindent Plusieurs ressources sémantiques ont été développées pour le domaine affectif. NRC Emotion Lexicon (Mohammad \& Turney, 2013) associe plus de 14\,000 mots à huit catégories d'émotions et deux polarités (positif/négatif) \cite{mohammad2013}. SentiWordNet étend WordNet par des scores de sentiment. EmoLex et EmoNet (Yang et al., 2007) fournissent des associations mots-émotions acquises par apprentissage sur de larges corpus.

\noindent Ces lexiques constituent une première source de connaissances utilisables par injection directe dans les représentations (\textit{feature concatenation}), par initialisation des embeddings ou par contraintes de régularisation sur les prédictions \cite{mohammad2013}.

\subsection{ConceptNet et le raisonnement sémantique}
\noindent ConceptNet (Speer et al., 2017) est un réseau sémantique multilingue à grande échelle, structurant des connaissances de bon sens sous forme de triplets (concept, relation, concept) avec des relations typées telles que \textit{IsA}, \textit{HasProperty}, \textit{Causes}, \textit{MotivatedByGoal}, ou \textit{CausesDesire} \cite{speer2017}. ConceptNet5 contient plus de 34 millions de triplets en plus de 300 langues \cite{speer2017}.

\noindent Son intégration dans les modèles de TALN s'effectue principalement via des \textit{Knowledge Graph Embeddings} (KGE) --- \textit{TransE}, \textit{RotatE}, \textit{ComplEx} --- qui projettent entités et relations dans un espace vectoriel continu tout en préservant les propriétés structurelles du graphe. Des approches récentes comme ATOMIC \cite{sap2019} enrichissent ConceptNet par des inférences causales et des chaînes de raisonnement événementiel, particulièrement pertinentes pour inférer les états mentaux implicites.

\subsection{Méthodes d'enrichissement des représentations par connaissances}
\noindent Plusieurs paradigmes d'intégration ont été proposés dans la littérature. La fusion par concatenation enrichit les embeddings contextuels des tokens par les embeddings de leurs voisins dans le graphe de connaissances. Les approches à base de graphes convolutionnels (GCN, GAT) propagent l'information sur le graphe de connaissances pour enrichir la représentation de chaque concept. ERNIE \cite{zhang2019ernie} réalise une fusion explicite entre les représentations BERT et les embeddings d'entités Wikidata via un mécanisme d'\textit{information fusion layer}.

\noindent Dans le contexte émotionnel, KnowBERT (Peters et al., 2019) et des approches comme KGNN (\textit{Knowledge Graph Neural Network for Emotion}) enrichissent les représentations des tokens ayant une valence affective en les ancrant dans leurs voisinages conceptuels dans \textit{ConceptNet} et \textit{EmoNet}, permettant ainsi d'inférer les émotions de tokens ambigus ou d'expressions idiomatiques.

\section{Techniques de Raisonnement par Bon Sens (Commonsense Reasoning)}

\subsection{Définition et importance pour la compréhension émotionnelle}
\noindent Le raisonnement par bon sens désigne la capacité d'un système à mobiliser des connaissances implicites sur le monde pour comprendre et inférer des informations non explicitement mentionnées dans un texte. Dans le contexte de la détection émotionnelle, cette capacité est cruciale : une phrase comme «~Il a raté son avion~» n'exprime pas directement une émotion, mais un système doté de connaissances de bon sens peut inférer la frustration ou l'inquiétude associées à cette situation.

\subsection{Ressources pour le raisonnement commun}
\noindent ATOMIC \cite{sap2019} est une base de connaissances d'inférences causales et sociales structurée autour de neuf dimensions relationnelles couvrant les causes (\textit{xNeed}, \textit{xWant}), les effets (\textit{xEffect}, \textit{xReact}, \textit{oReact}), et les réactions mentales. ATOMIC20 \cite{hwang2021} étend cette ressource à 1,33 million de tuples d'inférence et 23 relations.

\noindent COMET \cite{bosselut2019} est un modèle génératif entraîné sur ATOMIC capable de produire des inférences commonsense pour des phrases en entrée arbitraires, constituant ainsi un \textit{commonsense knowledge model} à la volée. COMET-ATOMIC 2020 combine les connaissances d'ATOMIC et de ConceptNet dans un modèle génératif unifié basé sur BART \cite{hwang2021}.

\subsection{Intégration du raisonnement commun dans les modèles émotionnels}
\noindent Des travaux récents ont exploré l'intégration du raisonnement commun pour améliorer la détection émotionnelle. Ghosal et al. (2020) proposent le modèle COSMIC qui enrichit un réseau récurrent par des inférences COMET pour chaque tour de dialogue, modélisant les états mentaux des interlocuteurs \cite{ghosal2020}. Ce modèle obtient des améliorations significatives sur les benchmarks de reconnaissance des émotions en conversation (ERC).

\noindent Ma et al. (2023) proposent une approche \textit{chain-of-thought} émotionnelle adaptant le paradigme de raisonnement en chaîne (Wei et al., 2022) au domaine affectif. En invitant un LLM à articuler une chaîne de raisonnement émotionnel explicite avant de produire sa prédiction, ils obtiennent des gains notables sur les tâches de détection émotionnelle \textit{zero-shot}. L'intégration de telles capacités de raisonnement explicite dans des architectures discriminatives (non génératives) reste un chantier ouvert.

\subsection{Inférence des émotions implicites}
\noindent Un sous-problème majeur de la détection émotionnelle est la capacité à identifier les émotions implicites, c'est-à-dire les émotions qui ne sont pas directement exprimées par des mots affectifs mais sont véhiculées par des descriptions d'événements, de situations ou d'actions. Des approches basées sur des modèles de causalité émotionnelle (\textit{emotion cause extraction}) et des réseaux bayésiens causaux ont été explorées. Des travaux comme ceux de Xia \& Ding (2019) sur l'extraction jointe causes-émotions montrent l'intérêt de modéliser explicitement le lien causal entre événements et réactions émotionnelles.

\section{Apprentissage Multi-Niveaux et Multi-Granularités}

\subsection{Granularités textuelles et représentations multi-échelles}
\noindent La modélisation des émotions à différentes granularités textuelles --- token, span, phrase, paragraphe, document --- constitue un axe de recherche actif. Les émotions peuvent être signalées par un mot isolé («~merveilleux~»), une expression figée («~avoir le cafard~»), une clause entière ou la tonalité globale d'un texte. Les approches à granularité unique, qu'elles opèrent au niveau du token ou du document, manquent nécessairement une partie de l'information émotionnelle.

\noindent Les architectures multi-résolution, inspirées des modèles de vision par ordinateur (\textit{Feature Pyramid Networks}), adaptent ce principe au texte en maintenant des représentations à plusieurs niveaux de granularité et en permettant des échanges d'information entre ces niveaux via des connexions latérales ou des mécanismes de fusion adaptative.

\subsection{Architectures multi-tâches pour l'apprentissage multi-niveaux}
\noindent L'apprentissage multi-tâches (\textit{Multi-Task Learning}, MTL) constitue un cadre naturel pour l'apprentissage multi-niveaux : différentes têtes de classification sont entraînées conjointement sur des tâches opérant à différentes granularités --- détection de spans affectifs au niveau token, classification émotionnelle au niveau phrase, et inférence du ton général au niveau document. Ruder (2017) et Liu et al. (2019) (MT-DNN) démontrent les bénéfices de la régularisation implicite apportée par le MTL pour les tâches de TALN \cite{liu2019mtdnn}.

\noindent Des architectures récentes comme ELECTRA fine-tuné en MTL sur des tâches émotionnelles, ou des modèles à encodeurs partagés avec des décodeurs spécialisés par granularité, illustrent cette tendance. La propagation de gradients à travers plusieurs niveaux d'abstraction favorise l'apprentissage de représentations plus généralisables.

\subsection{Pooling adaptatif et agrégation inter-niveaux}
\noindent La question de l'agrégation des représentations entre niveaux est cruciale. Des mécanismes de pooling appris (\textit{attention pooling}, \textit{gated pooling}) remplacent avantageusement le simple \textit{mean pooling} ou \textit{max pooling} pour agréger les représentations de tokens en représentations de phrases. De même, la représentation de document peut être construite par une attention hiérarchique sur les représentations de phrases, pondérant davantage les phrases émotionnellement saillantes.

\noindent Le modèle HAKE-MER du présent travail s'inscrit dans cette lignée en proposant des mécanismes de pondération dynamique adaptatifs permettant de gérer des émotions multiples d'intensités variables, surmontant ainsi la limitation des approches à seuil fixe qui traitent toutes les émotions de manière équivalente indépendamment de leur intensité.

\subsection{Évaluation multi-granulaire}
\noindent L'évaluation des modèles multi-étiquettes requiert des métriques spécifiques. Le F1-micro agrège les contributions de chaque instance uniformément, favorisant les étiquettes fréquentes. Le F1-macro calcule la moyenne non pondérée du F1 par étiquette, traitant équitablement toutes les émotions. L'\textit{Exact Match Accuracy} (ou \textit{Subset Accuracy}) mesure la proportion d'instances pour lesquelles l'ensemble exact d'étiquettes prédites coïncide avec l'ensemble de référence. Pour les applications sensibles à l'intensité, le \textit{Mean Average Precision} (mAP) et la corrélation de Pearson sur l'intensité émotionnelle (\textit{Emotional Intensity Correlation}) complètent ce tableau de bord.

\section{Synthèse et Positionnement de HAKE-MER}

\noindent Le tableau \ref{tab:synthese_hake_mer} suivant résume les forces et limitations des approches existantes au regard des objectifs du présent travail :

\begin{table}[H]
    \centering
    \caption{Synthèse des approches existantes et positionnement de HAKE-MER}
    \label{tab:synthese_hake_mer}
    \renewcommand{\arraystretch}{1.4}
    \setlength{\tabcolsep}{6pt}
    \begin{tabular}{|p{4.2cm}|p{4.2cm}|p{6.2cm}|}
        \hline
        \textbf{Approche} & \textbf{Forces} & \textbf{Limites} \\
        \hline
        
        \textbf{Méthodes classiques} \newline (SVM multi-label, Classifier Chains)
        & 
        Robustesse sur petits datasets, bonne interprétabilité
        & 
        Incapacité à capturer les dépendances contextuelles à longue portée et les émotions implicites
        \\
        \hline
        
        \textbf{Modèles BiLSTM avec attention}
        & 
        Modélisation séquentielle correcte, meilleure prise en compte du contexte que les approches classiques
        & 
        Sensibilité à la longueur des textes et absence de raisonnement sémantique
        \\
        \hline
        
        \textbf{BERT/RoBERTa fine-tuné}
        & 
        Représentations contextuelles riches et bonnes performances globales
        & 
        Traitement du texte à une granularité unique et absence d’intégration explicite de connaissances externes
        \\
        \hline
        
        \textbf{Approches avec connaissances externes} \newline (ERNIE, KnowBERT)
        & 
        Enrichissement sémantique et meilleure prise en compte des connaissances du monde
        & 
        Généralement mono-granulaire et sans mécanisme explicite d’interaction inter-émotions
        \\
        \hline
    \end{tabular}
\end{table}
\newpage 
\noindent HAKE-MER se distingue en combinant ces dimensions complémentaires dans une architecture unifiée : l'encodage hiérarchique multi-niveaux (tokens $\rightarrow$ phrases $\rightarrow$ document), l'attention croisée inter-émotions pour modéliser les corrélations et incompatibilités entre étiquettes, l'enrichissement sémantique par \textit{ConceptNet} et \textit{EmoNet} pour l'inférence des émotions implicites, et la pondération dynamique pour gérer les intensités différentielles. Cette combinaison vise une amélioration de 3 à 5\% en F1-score par rapport à l'état de l'art sur \textit{GoEmotions} et \textit{SemEval-2018}.




\section*{Conclusion}
\noindent Cet état de l'art a permis de baliser les grands axes de recherche constituant le socle théorique et méthodologique du présent PFE. La détection des émotions multiples dans les textes se trouve à la confluence de plusieurs champs : la classification multi-étiquettes, les architectures profondes à attention hiérarchique, l'ingénierie des connaissances et le raisonnement commun. L'architecture HAKE-MER s'appuie sur les avancées récentes de chacun de ces domaines pour proposer une approche intégrée, dont le plan expérimental --- basé principalement sur le dataset \textit{GoEmotions} avec validation croisée sur \textit{SemEval-2018} et \textit{EmoInt} --- permettra d'évaluer rigoureusement les contributions individuelles de chaque composante (étude d'ablation) et la robustesse de la généralisation inter-datasets.



